\section{Causal fragments and original cuts}\label{sec:fragments}

\begin{lemma}[Fragment starts]\label{lem:fragments}
\formal{Fragments.lean}{fragment_count_le,localized_fragment_count_le}
If a finite retained set $K$ is closed under state predecessors except at a
cut $B$, it has at most $|B|+1$ nonempty runs. If closure is required only
inside a selected set $C$ and $K\subseteq C$, it has at most
$|B|+r(C)$ runs, where $r(C)$ is the number of runs of $C$.
\end{lemma}
\begin{outline}
A noninitial run start must immediately follow a cut. In the localized
case it may instead be a run start of $C$. Map the other starts injectively
to their preceding cut times.
\end{outline}

\begin{lemma}[Current versions suffice]\label{lem:current-version}
\formal{CurrentVersion.lean}{filtered_read_correct,ExecutionPrefix.filtered_native_read}
Process an actual native execution's chosen writes in chronological order,
starting with blanks. At a requested read the resulting dictionary gives
the actual symbol if the actual last prior visit to that cell is processed,
or if no prior visit exists. No equality on other cells is required.
\end{lemma}
\begin{outline}
The recursive dictionary is unchanged on an interval with no write to the
queried cell. If a last visit exists, both the full and filtered dictionaries
therefore equal its written value. Otherwise both equal the initial value.
The native statement follows by induction from \cref{lem:old-write}.
\end{outline}

Let $\beta=q+ka+2$ encode the complete step value: poststate, all work writes
and optional output bit. Old positions and move trits are paid separately.
For $s\ge0$ put $n=s+1$ and
\begin{equation}\label{eq:global-menu}
 \mathcal J_s=(T_*+1)\binom{T_*}{s}\binom{T_*+n}{n}
 \left[6^s\binom{T_*+s}{s}\right]^{k+1}
 \binom{T_*+2n}{2n}\binom{P+4s+2}{4s+2}.
\end{equation}
The exact menu in \cref{eq:global-menu} is checked by
\lean{globalGeometryMenu_card}.
The declarations \lean{sampledWalk} and \lean{headDescription_injective}
construct the signed-increment description from chronological samples of a
unit walk. Monotone counters are encoded by their multiset, whose sorted
order recovers every sample. The actual native encoder and its raw-coordinate
and counter recovery are checked in \leanfile{FragmentGeometry.lean};
\lean{totalFragmentRecord_card} in \leanfile{FragmentMenuCount.lean}
combines them with the literal cut set and complete cut payload.

\begin{theorem}[Finite adaptive input-cut bound]\label{thm:input-cut}
\formal{NativeInputCutEndpoint.lean}{mem_nativeBadInputCuts_iff,native_bad_input_cut_bound}
For any fixed output permutation $\pi$, let $\operatorname{BadInput}(D,r)$
be the timely correct inputs admitting sets $V\subseteq U$ with $|V|\ge D$
and a vertex cut $B$, $|B|\le r$, from $\operatorname{Out}_X(V)$ to all of
$\operatorname{Read}_X(U)$. Then
\begin{equation}\label{eq:input-cut-count}
 |\operatorname{BadInput}(D,r)|\le
 2^{P-D}\sum_{s=0}^{\min(r,T_*)}\mathcal J_s2^{\beta s}.
\end{equation}
The sets and cut may depend on the input and are not separately encoded.
\end{theorem}
\begin{outline}
Retain the ancestors $K$ of the noncut desired roots in the graph with $B$
deleted. In each $B$-free gap the state spine makes $K$ a prefix. Every work
parent of a retained step belongs to $K\cup B$, unless the cell is virgin.
Record every cut value, old position and move, and the retained prefix lengths.
At a retained interval entry the state and heads are initial, or are obtained
from the immediately preceding cut step. Process retained transitions and
recorded cut writes in time order. By \cref{lem:current-version}, an omitted
overwrite cannot remain the latest parent of a retained read; closure would
have retained or cut it. A simultaneous induction compares two actual executions with the same
record. Their retained instructions agree because their old tuples agree;
no common detailed trace is assumed.

Record raw extrema for each retained interval and the before/after output
counters for every retained interval and every cut step. Let $I$ be the raw
coordinates read in $K$, and $O$ \emph{all} output ordinals at $K\cup B$.
The paid record determines both, including append flags at cut steps
through their two counters. Record the exterior of $W=\pi(O)\setminus I$. Retained
raw reads are known. Cut outputs are paid directly. Replay recovers the
outputs $O$ and therefore $W$, hence the entire input. The separator gives
$I\cap U=\varnothing$ and $V\subseteq\pi(O)$, so $V\subseteq W$.
Full correctness at all of $O$ is essential to this description-free step.
The formal event is expanded into an actual first halt, the complete output,
and the literal endpoint-inclusive separator by
\lean{mem_nativeBadInputCuts_iff}.

Here is the exact menu count. The horizon and cut set cost
$(T_*+1)\binom{T_*}{s}$. The $n$ prefix lengths have total at most $T_*$,
giving $\binom{T_*+n}{n}$. For each head, successive sampled cut positions
have total absolute variation at most $T_*$, giving
$2^s\binom{T_*+s}{s}$ signed-increment descriptions; moves add $3^s$.
The $2n$ nonnegative raw-extremum offsets have total at most $T_*$, giving
$\binom{T_*+2n}{2n}$. Finally the $4s+2$ chronological output-counter
samples are weakly increasing in $[0,P]$, giving
$\binom{P+4s+2}{4s+2}$. Empty intervals and invalid descriptions only
overcount. These factors are \cref{eq:global-menu}. The cut payload costs
$\beta s$ bits, and the exterior at most $P-D$ bits. Replay proves the
injection separately for every actual execution sharing the message; it
does not fix one input's query set for all other inputs. Sum over $s$.
\end{outline}

\begin{corollary}[Theorem I: simultaneous source-input expansion]\label{cor:input-expansion}
\formal{InputCutTail.lean}{native_input_cut_tail,native_input_cut_tail_all_width}
Fix $k,\beta$, a runtime constant $C\ge0$, and $\alpha>0$. There is
$\lambda>0$ such that at all sufficiently large lengths fewer than
$2^{(1-\alpha/3)P}$ timely correct inputs admit any cut in
\cref{thm:input-cut} with $|V|\ge\alpha P$ and $|B|\le\lambda P$.
For transposition the threshold is uniform over all $b\ge1$ as $d\to\infty$.
\end{corollary}
\begin{outline}
Put $r=\lfloor\lambda P\rfloor$. Bound every binomial factor in the exact
menu by the common sparse sum with $4r+2$ chosen positions out of
$P+T_*+4r+2$. The full cut menu is at most
\[
 (r+1)(T_*+1)\bigl[6^{k+1}2^\beta\bigr]^r
 \left[\sum_{j=0}^{4r+2}\binom{P+T_*+4r+2}{j}\right]^{k+5}.
\]
The generating-function sparse estimate and the polynomial logarithm bound
make its base-two logarithm at most $\alpha P/2$ after choosing one positive
$\lambda$ and then a length threshold. This is checked by
\lean{geometry_cut_menus_logb_small} in \leanfile{GeometryMenuRates.lean};
it requires no $r\le T_*/2$ restriction. The finite native bound gives
$2^{(1-\alpha/2)P}<2^{(1-\alpha/3)P}$ for $P>0$. The proof uses
$T_*\le\max(1,\lceil C\rceil)(P+1)$ and $P=d^2b\ge d^2$, so one threshold
works for every positive block width. No union over $U,V$ is needed.
\end{outline}

\subsection{Localizing the same decoder}

For a selected slab chunk $C$ from \cref{lem:slab-geometry}, keep its geometry
and ordinary entry records. Put $U=\pi(O(C))\setminus I(C)$. Let $Y_H$ be
\emph{all earlier endpoints} of off-carrier returns of excursion greater
than $H$, including endpoints outside $C$. Suppose $B$ separates
$\operatorname{Out}(U)$ from $Y_H$ in the full original graph. Retain the
ancestors $K$ of the noncut roots in the graph induced on $C\setminus B$.
There are at most $N+|B|$ retained fragments. A parent in $C$ is in $K\cup B$;
any parent in $Y_H$ is in $B$, even when it is outside $C$, by the full
separator property.

Process ordinary snapshots as pre-step overlays, and retained/cut writes as
step events. At the outer start install the free-slab snapshot or blanks. At
each original $C$-run entry overlay the actual nonfree windows. Do not erase
dictionary cells outside these windows. At ties perform overlays before the
cut or retained write. Fragment entry states and heads come from an original
run entry or the preceding cut step; record their output counters separately.

For a retained read on the free tape, its last visit after the outer start
belongs to $C$, hence to $K\cup B$; an older version is covered by the initial
slab snapshot. On a nonfree tape, a last visit in the current $C$-run is in
$K\cup B$. An earlier last visit is covered by the entry window, or lies in
$Y_H$ by \cref{lem:first-read} and is therefore in $B$. A virgin cell is
blank. These cases supply the correct current version. Later true snapshots
cannot spoil it: after the last visit and before this read every actual
snapshot of that cell has the same value. This also handles cut steps before
the outer start, cut roots and snapshot/cut ties. It would be false to claim
that a whole core dictionary remains correct after skipping $C\setminus K$.

The raw exterior is determined by the original $C$, so no new extrema for
the smaller fragments are required. With $s=|B|$, their extra geometry has
the safe menu
\begin{equation}\label{eq:local-menu}
 \mathcal L_s=(T_*+1)\binom{T_*}{s}\binom{T_*+N+s}{N+s}
 \left[6^s\binom{T_*+s}{s}\right]^{k+1}
 \binom{P+4s+2N}{4s+2N}.
\end{equation}
Its exact finite product cardinality is
\lean{localGeometryMenu_card}. The actual native local record and its
recovery are checked in \leanfile{LocalGeometry.lean}.
The bound \lean{native_bad_slab_cut_bound} combines the derived native
injection with the original slab metadata.
The factors count cut times, the $N+s$ prefix lengths, sampled positions and
moves, and $2(N+s)+2s$ output-counter samples. The proof is the same joint
sparse count as above. Replacing the import menu in \cref{eq:finite-returns}
by $\sum_{s\le r}\mathcal L_s2^{\beta s}$ therefore bounds the inputs with
such a cut of size at most $r$.

\begin{theorem}[Theorem H: original off-carrier expansion]\label{thm:off-carrier}
\formal{OffCarrierExpansion.lean}{native_off_carrier_expansion}
For fixed $k,a,q,C$ there are $\varepsilon,\lambda,\eta>0$ and a threshold
uniform over all $b\ge1$ such that, outside fewer than $2^{(1-\eta)P}$
timely correct inputs, every choice of free head admits a demanded raw set
$U$ with $P/(8N)\le|U|\le P/2$ whose original minimum vertex cut to all
off-carrier $\lfloor\varepsilon P\rfloor$-long return lower endpoints is
greater than $\lambda P$. In fact every valid chunk of
\cref{lem:slab-geometry} has this property on the same good-input event.
\end{theorem}
\begin{outline}
The exhaustive current-version induction above gives the actual native
pairwise injection in \lean{NativeReplayClass.LocalCut.records_recover}
of \leanfile{NativeLocalCuts.lean}. Its last-parent and desired-output
coverage are derived from the original full separator and actual counters.
Inputs sharing computed slab metadata have the same omitted set, even when
they halt at different later times. Group those actual inputs into metadata
fibres, pay the ordinary snapshots once, and sum their local-record bounds.
No common detailed trace or import bound is imposed.

Take $H=\lfloor\varepsilon P\rfloor$ with the explicit
$\varepsilon$ of \cref{thm:returns}. The common sparse envelope from
\lean{geometry_cut_menus_logb_small} makes the exact local cut menu's
logarithm at most $P/(32N)$ after choosing one positive $\lambda$ and a
length threshold. The windows cost at most $P/(32N)$ apart from fixed
terms. The metadata, the $3Ndb$ term, and all fixed terms including $Nq$
each cost at most $P/(96N)$ beyond a single dimension threshold.
The margin $P/(4N)$ therefore leaves exponent at most $1-5/(32N)$,
strictly below $1-\eta$ for $\eta=1/(8N)$.
This is \lean{native_local_cut_tail} in \leanfile{LocalCutTail.lean}.

The one bad event ranges over every valid actual chunk, so it already
includes all free-head and chunk choices. At the enlarged threshold the
geometric saving is at least $P/(8N)>0$; the proved slab-selection theorem
therefore supplies a nonempty valid chunk for every specified free head.
Its margin gives $|U|\ge P/(8N)$, and injectivity of actual append times
gives $|U|\le|O(C)|\le|C|\le P/2$. Decreasing $H$ only adds original targets,
so the same $U$ works for every smaller nonnegative integer threshold.
All length thresholds use $P=d^2b\ge d$, uniformly for $b\ge1$.
\end{outline}

The cut-density constant depends on the whole value width $\beta$, hence on
the controller. Unlike \cref{thm:returns}, this fixed positive cut-density
consequence does not automatically extend to arbitrary $o(P)$-bit controllers.
No result here supplies a small separator for every linear-time execution.
